\chapter{物质重构的物理疆界：基本力约束与广义黏滞模型}

第一章最后停在了一个更早的问题上：在一个由能量梯度驱动、由耗散维持的世界里，智能这种东西为什么会出现。这个问题不回答，后面所有关于资源、时间和结构的约束都只能被读成外部障碍，仿佛智能本来可以轻盈地穿过物质世界，只是碰巧遇到了阻力。真实的情形正好相反：阻力不是智能的外在障碍，它来自智能的起源方式。

智能不是从虚无中降临的奇迹，也不是生命偶然获得的一种装饰性能力，它更像是热力学世界在长期演化中生成的一种路径选择机制。要理解智能为什么出现，以及它为什么会遇到资源、时间和结构的约束，不能只从大脑、语言或计算机开始，而要从更深的地方开始：从能量梯度、耗散结构、自催化网络、实体试错和虚拟推演开始。

\section{智能与物质重构：优化算法的物理局限}

\subsection{最小作用量原理与路径选择器}

\subhead{能量梯度中的耗散结构}

宇宙并不静止。只要存在温差、势差、浓度差和化学能差，物质就会流动、碰撞、重组和耗散。孤立系统趋向熵增，但开放系统可以在持续能量和物质通量中形成局部有序结构。风暴、细胞、生命体、城市和机器，都不是熵增的反面，而是在更大范围熵增中形成的局部通道。它们维持秩序，不是因为逃离了热力学，而是因为吸收自由能、重组物质、排出废热，并以自身结构参与更大的能量释放过程。

这类局部秩序靠什么维持，需要一个更明确的机制说明。薛定谔在《生命是什么》中提出一个影响深远的问题：生命体如何避免迅速滑向热平衡？他的回答是，生命通过从环境中摄取“负熵”来维持自身秩序。\cite{schrodinger1944} 用今天更稳妥的说法，生命不是违背热力学第二定律，而是开放系统在持续输入低熵资源、输出高熵废物和废热的过程中维持局部有序。生命的秩序不是免费获得的，它必须不断支付代价。

Prigogine 的\term{耗散结构}理论把这个问题进一步推进。远离平衡态的开放系统，在持续通量和非线性相互作用下，可能形成新的有序结构。\cite{prigogine1977nobel,nicolis1977} 这种有序不是静止晶体式的封闭秩序，而是依赖输入和输出的动态秩序。它只要停止获得能量和物质流，就会瓦解、冷却、死亡或回到更接近平衡的状态。生命正是这种动态有序的极端形式之一：它不是热力学之外的例外，而是热力学内部一种复杂、持久、可复制、可演化的耗散结构。最大熵产生原理可以作为一种启发性图像：在一些远离平衡的开放系统中，能够稳定组织通量的结构，可能更容易被维持和放大。\cite{martyushev2006} 风暴组织热量与动量流，生物体捕获并转化化学能或太阳能，生态系统和城市则在更大尺度上重组能量、物质和信息流。

但这里必须谨慎：最大熵产生原理不是一个无争议、可直接套用到所有生命和智能演化的普遍定律。\cite{bruers2007} 本章只把它作为图像，而不是证明。它帮助我们看见一种可能：\key{秩序并不必然意味着对熵增的抵抗；很多时候，秩序本身就是更复杂、更稳定、更持久的耗散通道。}

\subhead{自催化与条件的耦合}

局部有序一旦形成，接下来要解决的是延续问题。从普通化学反应到生命，中间需要一些关键过渡，其中最重要的一类，是自催化。在原始海洋、热泉、矿物表面或小尺度化学环境中，分子会因为热运动、随机涨落和局部浓度差形成短暂的小序。大多数小序很快解散，重新回到无序分布。但总会有一些结构比较特殊：它们不仅存在，而且会促进自身形成。A 和 B 反应生成 C，而 C 又促进 A 和 B 继续反应，这就是自催化。自催化把偶然出现的小序放大，使短暂结构获得自我强化的能力。

自催化网络之所以重要，是因为它把“偶然结构”变成了“可延续结构”。它不只是一次化学反应，而是一个正反馈网络。只要有持续的能量和原料输入，这种网络就可能扩展、复制、竞争和分化。它吸收自由能，维持局部秩序，同时向外排放热量和废物。更复杂的结构能够更有效处理物质流和能量流，也更容易在竞争中留下来。Kauffman 关于自组织和自催化集合的工作，正是试图说明生命起源不能只从外部选择来理解，还要理解化学网络内部如何自发形成可维持、可扩展的结构。\cite{kauffman1993}

这里也不能把问题说得过满。自催化网络不是生命起源的完整答案。Hordijk 等人讨论自催化集合时也强调，自催化集合是类生命行为的重要条件，但并不足以单独构成完整生命。\cite{hordijk2018} 一个自我强化的反应网络，还需要边界、信息保存、变异、复制、选择和内部调节，才可能接近生命。自催化提供的是从偶然小序走向可延续结构的一条路径，而不是把生命起源一次性解释完毕。

缺的那几项条件各有来源，其中被讨论最多的是信息载体。RNA 世界假说在这里提供了另一条重要线索。RNA 的特殊性在于，它既可以承载序列信息，又在某些情形下具有催化功能。Walter Gilbert 在 1986 年用“RNA world”这一说法描述一种假想阶段：在 DNA 和蛋白质主导的现代生命之前，RNA 可能曾同时承担信息保存和催化功能。\cite{gilbert1986} 这个假说并不等于生命起源已经被解决，但它说明一个关键点：生命的出现需要把信息、反应和复制逐渐扣在一起。只有能量流，没有信息，结构很难积累历史；只有信息，没有反应和能量，结构无法维持；只有反应，没有边界，局部秩序又难以集中和继承。

因此，耗散结构、自催化网络、RNA、脂质膜和原始代谢，不应被理解为一条已经被完全证明的单线历史，而是几类必须逐渐耦合的条件：能量通量、边界、信息保存、复制、变异、选择和内部调节。膜把系统与环境区分开来，RNA 或其他信息载体保存可继承的结构，自催化网络把能量流和物质流组织起来。原始生命不是突然出现的灵魂，而是某些耗散结构在长期选择中获得边界、复制和内部调节能力的结果。

\subhead{相空间中的测地线}

耗散结构和自催化网络说明了局部有序如何出现、如何延续，但没有说明智能在这套图景中占据什么位置。要回答这个问题，需要换一种描述方式：把一次制造活动看成系统状态在一个高维空间中的移动。

日常所说的空间只有长、宽、高三个维度，一个单摆的状态只需要位置和动量两个维度就能确定。而一座化工厂或一条总装线的状态，要用成千上万个独立变量才说得清：反应温度、压力、催化剂浓度、材料强度、物流距离、库存周期、设备折旧率、资金成本，每一项都是一个维度。原材料是其中一个坐标点，成品是另一个坐标点，生产过程就是状态点在两者之间画出的轨迹。所谓工艺改进，在这个描述里就是轨迹的重新选择。

在状态空间里，工程优化是在约束下比较路径：同一个目标，能否用更少材料、更低能耗、更短工期或更小风险达到。物理学中的驻作用量原理提供另一种严格的路径表述。它从规定的拉格朗日量与边界条件导出运动方程，不能直接等同于经济成本最小化。两种表述都关注路径，却采用不同的量与判据。

\begin{equation}
S[q]=\int_{t_{1}}^{t_{2}}\left(T-V\right)dt,
\qquad
\delta S=0
\end{equation}

$S$为作用量，单位为焦耳秒；$q$为广义坐标，$T$与$V$分别为动能和势能；$\delta S=0$表示真实轨迹在规定端点条件下使作用量取驻值。驻值可以是极小值，也可以有其他性质。摩擦等非保守作用需要相应的广义力或耗散描述。这个原理提供运动方程，智能如何形成选择、记忆和学习，还需要另作机制说明。

工业中的浪费有具体来源：需求预测不准形成库存，工况控制不足增加损耗，路线安排不良造成重复运输。AI可以通过搜索、建模与控制减少这些支出。改进的大小，应在相同任务、质量和系统边界下测量，不能用“额外作用量”代替能耗、费用与时间各自的账本。

智能在这里是一种\term{路径选择器}。它在目标、可用资源与反馈条件下比较行动，寻找成功率更高、代价更可承受的路线。这个选择可以改变材料、工艺与初始状态，也可以改变实际损失；具体下限则要随任务和物理边界重新计算。

路径选择把一部分昂贵试错提前放进模型。比较本身有成本，所选路径也要在现实中执行。接下来需要分清三笔账：反应的热力学净代价、反应速率所受的势垒，以及设备和结构的运行维护。三者相互影响，却不是同一笔能量。

\subsection{两笔刚性支出：活化能与熵税}

改变与维持受不同条件约束。活化势垒主要影响反应速率，自由能差约束给定条件下的净功交换，维护则涉及损伤、扰动与功能持续。晶体等稳定结构可以被动维持，生命与持续运行的设备通常需要通量支持。把它们分开，才看得清智能能改进哪一段。

\subhead{实体试错的热力学瓶颈}

但生命一旦出现，就面对一个更严酷的问题：如何在复杂环境中找到更好的路径。最初的生命系统主要依赖\term{实体试错}。突变、复制、竞争、淘汰和死亡，构成最古老的搜索机制。它有效，因为现实会筛选；但它极其昂贵，因为每一次失败都在真实世界中结算。一次错误突变不是修改一行代码，而是改变一个身体；一次错误行动不是改正一个答案，而是消耗能量、暴露风险、丧失繁殖机会，甚至死亡。自然选择的计算发生在真实环境中，它的代价由物质、时间和生命支付。

这就是实体试错的热力学瓶颈。环境越复杂，可能路径越多，纯粹依靠现实碰撞来筛选路径就越昂贵。演化可以用海量个体和漫长时间进行搜索，但对每个具体生命体来说，现实试错太慢、太重、太危险。一个个体不能每次都用身体去验证悬崖能不能跳、火能不能靠近、陌生动物会不会攻击、腐坏食物能不能吃。现实会给出答案，但答案常常来得太晚。

智能正是在这个瓶颈中出现的。神经系统的意义，不只是让生命感知世界，而是在身体内部建立一个低成本的环境模型。动物不必把每一种可能行动都拿到现实中试一遍，而可以先在内部预演：哪里可能有食物，哪里可能有危险，哪条路线更省力，哪个动作更可能成功，哪个对象值得接近，哪个信号必须回避。智能把一部分外部碰撞转移到内部，把一部分实体试错转化为\term{虚拟推演}。

这是一场深刻的成本置换。现实中的试错需要位移、受伤、能量消耗、材料损耗甚至死亡；内部推演需要神经信号、记忆、注意力和计算。二者都不是免费的，但数量级不同。智能使系统能够在做功之前先选择，在行动之前先比较，在暴露于风险之前先排除错误。它不取消物理代价，而是减少盲目支付物理代价的次数。

从生命演化和行动系统的角度看，\key{智能的第一性功能不是抽象思考，而是少犯昂贵的错误。}这个判断听起来不像传统哲学中对智能的赞美，却更接近生命的处境。一个生命体必须在有限能量、有限时间、有限身体和有限信息中行动。它不需要知道宇宙真理才算智能；它首先需要在行动前压缩风险，在接触前判断对象，在消耗前选择路径，在失败前排除明显错误。智能不是抽象心灵的豪华能力，而是身体在沉重世界中延续自身的路径选择机制。

\subhead{信息中枢与两类熵}

这个机制一旦运转，就会遇到一个更细的问题：它压缩的究竟是什么。智能通过内部模型降低主体对相关环境状态和行动后果的不确定性，也就是降低行动者内部的信息熵。但这里需要区分两类熵。热力学熵描述能量耗散和不可逆过程；信息熵描述系统对状态的不确定性。Shannon 建立信息论时，把不确定性形式化为信息熵。\cite{shannon1948}

智能降低的不是外部世界本身的热力学熵，而是主体对环境和行动后果的无知程度，使系统更清楚地判断路径；而这种信息熵的降低，必须以真实物理过程为载体。大脑需要葡萄糖和氧气，神经信号需要电化学过程，记忆和学习需要物质结构变化。智能不是无成本的精神活动，而是开放系统内部一种高度组织化的耗散过程。

人脑本身就是一个清楚的事实锚点。Raichle 和 Gusnard 在讨论脑能量预算时指出，人脑消耗约占身体氧气和热量消耗的 20\%，尽管它只占身体质量的一小部分。\cite{raichle2002} 这个比例本身就说明，智能不是免费的。大脑之所以昂贵，是因为它把大量环境信息、身体状态、记忆、预测和行动选择压缩在一个高耗能系统中。生命愿意为大脑支付高代价，恰恰说明内部推演能够节省更昂贵的外部试错。

兰道尔原理（Landauer）也提醒我们，信息处理不能脱离物理成本。他在 1961 年指出，逻辑不可逆操作中的信息擦除会伴随热量产生。\cite{landauer1961} 这个原则不能被粗暴地用来计算所有思考或所有 AI 训练的能耗，但它足以说明一个基本事实：信息不是脱离物质世界的幽灵。信息的记录、擦除、传输和计算都需要物理基座。智能降低不确定性，但不脱离热力学；智能节省错误行动的代价，但自身仍然要消耗能量。

生命体内部也因此逐渐发生功能分化。一个复杂生命系统需要能量系统、信息系统和做功系统。能量系统负责摄取、转化和供应自由能；做功系统负责移动、捕食、建造、逃避和改变环境；信息系统负责感知、记忆、预测和选择路径。细胞中有能量代谢、遗传信息和蛋白质机器；多细胞生物中有消化、循环、神经和运动系统；人类文明中则有能源基础设施、生产体系、知识体系和计算系统。智能就是耗散结构中的信息中枢。它自己不直接替代能量，也不直接替代做功，但它决定能量和做功投向哪里。大脑之所以昂贵，是因为它压缩了大量环境信息；文明之所以发展出语言、文字、数学和科学，是因为这些系统能把分散经验压缩成可传递、可推演、可复用的路径选择能力。信息越复杂，压缩越高级，系统越能减少盲目试错。

Friston 的自由能原则与这个方向有某种相通之处。自由能原则试图把感知、行动和学习放进一个统一框架：有机体通过降低预测误差或惊讶来维持自身状态。\cite{friston2010} 本章并不把智能简单等同于自由能最小化，也不把一个宏大理论直接套进所有智能现象。这里要借用的只是其中一个深层直觉：生命不是被动接收世界，而是不断预测、行动、修正，以维持自身在可生存范围内。智能不是在世界之外旁观，而是在身体和环境的闭环中不断选择路径。

\subhead{启动势垒：断裂化学键的代价}

物质重构要经过实际反应路径。反应物与产物的自由能差，决定给定条件下净变化的热力学方向；活化自由能则影响跨越过渡态的速率。断开旧键需要能量，形成新键也会释放能量，净账不能只加断键支出。制造并不总把物质推向更高能状态，工艺的意义在于让目标状态以可行速率、纯度与收率出现。

在恒温恒压及相应反应条件下，吉布斯自由能变化可写成$\Delta G=\Delta H-T\Delta S$。负的反应自由能变化说明给定状态下的自发方向，却不提供速率。反应商、压力、温度与组分改变以后，方向与平衡位置也会变化。催化剂改变路径和速率，不改变同一初末状态的自由能差。

炼铁要把氧从含铁原料转移到其他物质中。高炉中的碳和一氧化碳参与还原，焦炭还提供热量并支撑料柱；能耗取决于原料、还原路线、温度、传热与设备。\cite{babich2021blast} 算法可以改善配料、流动和控制，替代还原剂与工艺路线也可能改变整套支出。单个铁氧键的键能不能直接当作每吨铁的最低能源账单。

水泥熟料生产包含碳酸盐分解与后续矿物形成。碳酸钙分解吸热，平衡还受二氧化碳分压影响；熟料烧成则需要规定的温度与停留条件。文献给出的约$1.7\,\mathrm{GJ/t}$理论热耗，是特定熟料组成与边界下的过程估算，实际热耗还包括传热、废气与设备损失。\cite{madivate1999theoretical,madlool2011cement} 改变熟料比例、原料或胶凝材料路线，比较基准也随之变化。

高温同时影响平衡、反应速率与传递过程。提高温度能加快某些步骤，也可能增加损失或损害产物；催化与路线重设计则可能在较低温度下实现同一目标。需要区分的是给定初末状态的热力学约束、所选路径的动力学障碍和实际设备的损耗。智能可以帮助改变后两者，也可以帮助寻找新的初始状态与目标组合。

\subhead{熵税：维持秩序与对抗随机性}

“熵税”在这里指维持功能所需的持续支出。生命依靠代谢、修复与调节，运行中的设备需要供能、散热与保养。建筑、晶体与被动储存结构则有不同的稳定机制，并非每一份有序都需要持续外部做功。维护成本应从具体损伤和扰动机制估计。

农业把光、水、营养、土地与生物过程组织成可收获的产出。育种可以改变收获指数，也可以改善抗病、光合、根系与资源利用。产量和投入不是固定总额下的单纯重新分配：它们会随品种、环境与管理共同改变。应当分别核算太阳能通量、外部投入、土地占用和收获结果。

建筑安全需要处理荷载、材料、施工与模型的不确定性。安全系数和分项系数按规范、失效模式与构件确定，不能统一解释成“理想一吨、现实两吨”。冗余材料也可能参与承载与耗能。智能设计的目标是满足可靠度要求下减少不必要支出，同时保留经得起扰动的结构。

兰道尔界限给出一种明确的微观账本：在相应热环境与信息条件下，将一个无偏比特复位，最低排热为$k_{\mathrm B}T\ln2$。室温下约为$2.9\times10^{-21}$焦耳。把一千克物体抬升一米约需9.8焦耳，两者选择了完全不同的操作单位；这个比值不能估算完整推理与完整制造任务的成本。比较完整任务，还要计入操作次数、实际器件、通信、设备及辅助系统。

\subsection{物质重构的守恒方程}

\subhead{虚拟推演的放大}

人类文明延续并放大了这条路线。语言把个体经验压缩成群体记忆，地图把空间试错提前为路径表示，数学把关系压缩成可推演形式，科学理论把分散经验组织成可检验模型，工程仿真把昂贵实验提前到表示空间。文明的许多发明，本质上都是把实体试错提前为符号试错、模型试错和计算试错。人不必真的走遍每条路，地图可以提前筛掉坏路径；工程师不必让每座桥都塌一次，力学模型可以提前筛掉危险结构；医生不必把所有治疗方案直接施加到病人身上，实验、动物模型、临床试验和统计分析共同压缩错误空间。

人工智能进一步提高了虚拟推演的密度。它能在巨大的可能空间中搜索候选方案，发现隐含结构，生成模型，预测结果，并在真实行动发生之前进行大规模筛选。AI 不是脱离热力学的超物理智能，而是这条演化路线的新阶段。大脑用神经信号模拟世界，科学用方程和实验体系模拟世界，工程用计算机仿真模拟世界，AI 用高维统计模型和生成系统模拟世界。它们的共同功能都是：在真正支付物理代价之前，先压缩错误、筛选路径、降低盲目性。

虚拟推演把选择和执行分开了，这个分开之处可以用一个经典的思想实验说明。麦克斯韦设想有一个守在容器阀门边的小妖，能看清每个分子的速度，只放快分子通过、拦住慢分子，从而在不消耗能量的情况下制造温差。Leo Szilard 和兰道尔后来证明，妖的测量和记忆要付出信息代价，这部分上一节已经算过；更关键的是，妖没有手。它能决定谁通过，却不能推动一个分子，更不能推动一扇有质量的阀门。工业系统里的人工智能处在这个小妖的位置上：它拥有极低的信息熵，可以看清高炉里的气流涡旋、物流网络里的车辆轨迹、电网里的频率波动；它不产生力，所有对原子状态的改变都要通过执行器完成，而执行器的运转依靠电流、液压和燃烧。从控制论的角度看，工业文明是一条功率放大链：纳瓦级的神经电信号控制瓦级的肌肉，瓦级的肌肉控制千瓦级的汽车，千瓦级的汽车控制兆瓦级的工厂。智能位于这条链的最顶端，也就是能级最低的信息控制层。控制信号不等于功率流，一道逻辑门可以决定水坝闸门是开还是关，却提供不了提起闸门所需的机械能；能源不足时，无论指令多么正确、多么及时，闸门都不会动。这条链条的功率关系可以写成连乘，

\begin{equation}
P_{k}=P_{0}\prod_{i=1}^{k}G_{i},
\qquad
G_{i}=\frac{P_{i}}{P_{i-1}}\gg 1
\end{equation}

式中 $P_{0}$ 为链条起点的控制信号功率，单位瓦；$P_{k}$ 为末级执行功率，单位瓦；$G_{i}$ 为第 $i$ 级的功率增益，无量纲；$k$ 为级数，正文给出的链条是纳瓦级的神经电信号、瓦级的肌肉、千瓦级的汽车、兆瓦级的工厂。每一级都不把上一级的信号当作能源，只把它当作开关，功率由本级自带的能源系统提供，所以整条链的功率从纳瓦级延伸到兆瓦级，靠逐级相乘完成。$G_{i}\gg1$ 还说明另一件事：信号层携带的功率与执行层相比可以忽略，因此“指令正确”与“能够动作”是两个彼此独立的条件。开关可以决定谁通过，却不能推动一个分子；上面这条连乘关系就是这个区别的算术形式。它的适用边界是：连乘只在每一级都配备独立功率源时才成立；任何一级缺少能量或执行机构，$G_{i}$ 降到 1 以下，链条就在该处中断，上游再高的计算精度也无法补偿。

\subhead{虚拟推演及其边界}

但虚拟推演不是现实本身。智能可以选择路径，不能替代路径；可以调度能量，不能取消能量；可以减少无效耗散，不能让物质重构不再耗散。它不是力，而是对力的安排；不是能量，而是能量使用方向的选择机制。智能的本质不是让世界变轻，而是在世界依然沉重的前提下，减少错误搬动世界的次数。

AI 可以设计一座更省材料的桥，却不能取消钢材的质量、重力势能和连接强度；可以优化机器人动作，却不能替电机输出力矩、替关节承受磨损、替传感器处理噪声；可以筛选药物候选，却不能替人体完成代谢、毒理、临床验证和长期反馈。\key{智能可以逼近物理效率的下限，却不能把物理下限归零。}它可以减少错误路径，不能取消路径本身；可以减少无效实验，不能取消验证；可以减少浪费，不能取消做功。

物质重构需要把描述、投入、过程与结果接起来，可以先写成一条核算关系：
\begin{equation}
\text{信息（配方）}+\text{物质（载体）}+\text{自由能（做功）}+\text{时间（过程）}\longrightarrow \text{有序结构}+\text{废热}
\label{eq:conservation}
\end{equation}
式中的四项输入各自对应一类不可替代的条件。信息是配方，量纲上表现为比特，它决定了结构该长成什么样，但本身不改变任何原子的位置；物质是载体，单位是千克，它提供最终结构所需要的原子；自由能是做功项，单位是焦耳，它提供把原子从低能态搬到高能态所需的能量；时间是过程项，单位是秒，因为状态转移必须沿一条真实的时间轨迹展开，不能只在表示空间中跳跃。右侧的两项是结果：有序结构是被固定下来的低熵状态，废热是过程无法避免的熵增输出。这个关系成立的条件是系统对外开放，允许自由能注入和废热排出；对一个封闭系统，右侧的废热项会吞掉全部输入，有序结构无法积累。它不是会计意义上的平衡式，而是过程意义上的单向条件：只有信息，没有物质，结果不会出现；只有物质，没有自由能，状态不会改变；只有能量，没有时间，过程不能展开；只有局部有序，没有外部熵增，系统不能维持。智能能改变的是选择，不是存在条件；能优化的是路径，不是现实的物理账本。
\begin{figure}[ht]
\centering
% >>>>>>>>>> figs/fig-virtual-trial.tex >>>>>>>>>>
% 图 2-2  虚拟推演把试错成本搬了家，没有把它取消
% (a) 双对数代价曲线：虚拟筛选与候选数成正比（斜率 1）；实体验证的单位成本
%     高两个量级（对数坐标里整条线上抬 2 格），斜率也更陡。两者相加给出
%     最优候选数 N*——再往上加候选，负担全部落到实体那一侧。曲线只画到
%     N=10^5，右半幅留白放图例，避免用引线把图面割碎。
% (b) 同一笔账按环节分摊：做虚拟推演后筛选段变厚，实体段只薄一点——
%     筛选能减少**试错次数**，不能减少**验证本身**，总高度几乎不变。
\begin{tikzpicture}[x=1cm, y=1cm, font=\small,
  ax/.style={-{Stealth[length=2.0mm]}, ink, line width=0.6pt},
  yk/.style={anchor=east, font=\scriptsize, text=slate},
  tk/.style={anchor=north, font=\scriptsize, text=slate},
  sw/.style={line width=1.6pt}
]
% 下边只剩 0.6pt 余量，(a)/(b) 小标题再往下一点点就会被切。放 0.16cm 保险。
\useasboundingbox (-0.95,1.00) rectangle (12.80,-7.06);

% ================= (a) 代价曲线 =================
\def\aX{1.05} \def\aY{-4.70} \def\u{0.62}
\draw[ax] (\aX,\aY) -- ({\aX+\u*7.55},\aY)
      node[anchor=west, font=\scriptsize, text=ink] {候选数 $N$};
\draw[ax] (\aX,\aY) -- (\aX,{\aY+\u*7.4})
      node[anchor=south, font=\scriptsize, text=ink, yshift=1pt] {累计代价};
\foreach \e in {0,1,2,3,4,5,6} {
  \pgfmathsetmacro\yy{\aY+\u*\e}
  \draw[slate!55, line width=0.4pt] (\aX,\yy) -- ({\aX-0.07},\yy);
  \node[yk] at ({\aX-0.09},\yy) {$10^{\e}$};
}
\foreach \e/\t in {0/1, 2.5/{$10^{2.5}$}, 5.0/{$10^{5}$}} {
  \pgfmathsetmacro\xx{\aX+\u*\e}
  \node[tk] at (\xx,{\aY-0.12}) {\t};
}

\draw[deepblue!85, line width=1.4pt] ({\aX},{\aY+\u*0.15}) -- ({\aX+\u*5.0},{\aY+\u*5.15});
\draw[rust!90, line width=1.4pt] ({\aX},{\aY+\u*2.0}) -- ({\aX+\u*3.0},{\aY+\u*6.5});
\draw[ink!80, line width=1.3pt, dashed] ({\aX},{\aY+\u*0.60})
      -- ({\aX+\u*1.30},{\aY+\u*1.86}) -- ({\aX+\u*3.9},{\aY+\u*5.20});

\pgfmathsetmacro\xstar{\aX+\u*1.30}
\draw[slate!60, line width=0.45pt, dotted] (\xstar,\aY) -- (\xstar,{\aY+\u*1.86});
\node[font=\scriptsize, text=ink, anchor=north] at (\xstar,{\aY-0.12}) {$N^{*}$};
\node[font=\scriptsize, text=slate, anchor=north west, align=left]
      at ({\aX+\u*0.15},{\aY-0.58}) {再往上加候选数，负担全落在实体那一侧};

% 图例块：放在曲线右侧空白，条目自带色线，不用引线
\def\gx{\aX+\u*5.55}
\draw[deepblue!85, sw] ({\gx},{\aY+\u*6.30}) -- ({\gx+0.55},{\aY+\u*6.30});
\node[font=\scriptsize, text=ink, anchor=west] at ({\gx+0.65},{\aY+\u*6.30})
      {虚拟筛选 $\propto N$};
\draw[rust!90, sw] ({\gx},{\aY+\u*5.50}) -- ({\gx+0.55},{\aY+\u*5.50});
\node[font=\scriptsize, text=ink, anchor=west] at ({\gx+0.65},{\aY+\u*5.50})
      {实体验证 $\propto N^{1.5}$};
\draw[ink!80, sw, dashed] ({\gx},{\aY+\u*4.70}) -- ({\gx+0.55},{\aY+\u*4.70});
\node[font=\scriptsize, text=ink, anchor=west] at ({\gx+0.65},{\aY+\u*4.70})
      {总代价};

% ================= (b) 环节分摊 =================
\def\bX{8.20} \def\bY{-4.70} \def\bw{1.00}
\draw[ax] (\bX,\bY) -- ({\bX+3.10},\bY);
\draw[ax] (\bX,\bY) -- (\bX,{\bY+3.85})
      node[anchor=south, font=\scriptsize, text=ink, yshift=1pt] {总代价};
\def\lX{\bX+0.40} \def\rX{\bX+1.75}
% 左柱：不做虚拟推演
\fill[deepblue!22] ({\lX},{\bY+0.02}) rectangle ({\lX+\bw},{\bY+0.62});
\draw[deepblue!70, line width=0.7pt] ({\lX},{\bY+0.02}) rectangle ({\lX+\bw},{\bY+0.62});
\fill[rust!30] ({\lX},{\bY+0.62}) rectangle ({\lX+\bw},{\bY+2.72});
\draw[rust!75, line width=0.7pt] ({\lX},{\bY+0.62}) rectangle ({\lX+\bw},{\bY+2.72});
\fill[terra!30] ({\lX},{\bY+2.72}) rectangle ({\lX+\bw},{\bY+3.45});
\draw[terra!80, line width=0.7pt] ({\lX},{\bY+2.72}) rectangle ({\lX+\bw},{\bY+3.45});
\node[font=\scriptsize, text=ink, anchor=north, align=center] at ({\lX+\bw/2},{\bY-0.14})
      {不做虚拟\\推演};
% 右柱：做虚拟推演
\fill[deepblue!22] ({\rX},{\bY+0.02}) rectangle ({\rX+\bw},{\bY+1.28});
\draw[deepblue!70, line width=0.7pt] ({\rX},{\bY+0.02}) rectangle ({\rX+\bw},{\bY+1.28});
\fill[rust!30] ({\rX},{\bY+1.28}) rectangle ({\rX+\bw},{\bY+2.92});
\draw[rust!75, line width=0.7pt] ({\rX},{\bY+1.28}) rectangle ({\rX+\bw},{\bY+2.92});
\fill[terra!30] ({\rX},{\bY+2.92}) rectangle ({\rX+\bw},{\bY+3.42});
\draw[terra!80, line width=0.7pt] ({\rX},{\bY+2.92}) rectangle ({\rX+\bw},{\bY+3.42});
\node[font=\scriptsize, text=ink, anchor=north, align=center] at ({\rX+\bw/2},{\bY-0.14})
      {做虚拟推演};
% 两柱之间：实体段厚度的对比括线
\draw[ink!60, line width=0.5pt, {Stealth[length=1.4mm]}-{Stealth[length=1.4mm]}]
      ({\bX+1.55},{\bY+0.62}) -- ({\bX+1.55},{\bY+2.72});
% 图例：贴在柱右侧，条目自带色块
\def\gX{\bX+2.98}
\draw[terra!80, line width=3.0pt] (\gX,{\bY+3.42}) -- (\gX+0.34,{\bY+3.42});
\node[font=\scriptsize, text=ink, anchor=west] at ({\gX+0.42},{\bY+3.42}) {返工};
\draw[rust!75, line width=3.0pt] (\gX,{\bY+2.60}) -- (\gX+0.34,{\bY+2.60});
\node[font=\scriptsize, text=ink, anchor=west] at ({\gX+0.42},{\bY+2.60}) {实体};
\draw[deepblue!70, line width=3.0pt] (\gX,{\bY+1.78}) -- (\gX+0.34,{\bY+1.78});
\node[font=\scriptsize, text=ink, anchor=west] at ({\gX+0.42},{\bY+1.78}) {虚拟};
% 括线注记放在柱顶上方，两条细引线分别指到两柱的实体段
\node[font=\scriptsize, text=ink, anchor=south] at ({\bX+1.55},{\bY+4.05})
      {实体那一段只薄了一点};
\draw[ink!45, line width=0.4pt] ({\bX+0.95},{\bY+4.02}) -- ({\bX+0.90},{\bY+2.00});
\draw[ink!45, line width=0.4pt] ({\bX+2.15},{\bY+4.02}) -- ({\bX+2.25},{\bY+2.30});

% ================= 栏标题（一律在内容下方）=================
\node[anchor=north, font=\footnotesize\heiti, text=slate] at (3.45,-6.55) {(a) 最优候选数};
\node[anchor=north, font=\footnotesize\heiti, text=slate] at (10.10,-6.55) {(b) 成本只是搬家};
\end{tikzpicture}

% <<<<<<<<<< figs/fig-virtual-trial.tex <<<<<<<<<<

\caption{虚拟推演没有取消试错成本，只是把它从实体一侧移到了虚拟一侧。(a) 虚拟筛选的代价与候选数成正比，实体验证的单位成本高两个量级且斜率更陡，两者相加给出最优候选数 $N^*$。(b) 把同一笔账按环节分摊开：两段的相对厚度变了，柱子的总高几乎没动。}
\label{fig:virtual-trial}
\end{figure}


\subhead{候选膨胀与代价重现}

这正是 AI 时代最容易被误解的地方。AI 降低的是找到可能路径的成本，而不是自动降低完成状态转移的全部成本。它可以生成一千种建筑方案，但建筑仍要施工；可以提出一万个分子候选，但实验和临床不能同步扩张；可以模拟无数机器人动作，但真实机器人仍要面对摩擦、碰撞、电池、关节和维护。搜索空间可以快速膨胀，验证和执行仍要回到物理世界。

更进一步，AI 本身也有热力学成本。模型训练和推理需要芯片、电力、数据中心、冷却系统、供应链和维护体系。所谓云端并不轻盈，它只是把材料、能耗和废热隐藏在界面之后。AI 内部表现为信息压缩和路径搜索，外部表现为高密度耗散结构。它把一部分实体试错转化为计算试错，但计算试错仍然需要物理基座。用比特替代一部分原子试错，并不等于没有原子；用计算替代一部分实验，并不等于没有能耗；用模型筛选路径，并不等于路径已经完成。

智能越强，现实越不会自动变轻。智能首先扩展的是候选空间。过去只能尝试十种方案，现在可以生成一万种；过去只能模拟少数条件，现在可以探索巨大参数空间；过去只能依赖少数专家直觉，现在可以让模型压缩海量经验。这当然是进步。但候选变多，也会把压力推向验证、执行和承受。智能降低了搜索成本，却可能提高筛选负担；降低了提出方案的门槛，却提高了确认方案的要求；降低了比特空间中的试错成本，却把原子和生命空间中的验证瓶颈显影出来。

智能形成内部表示，使一部分昂贵试错可以先在模型中进行。技术黏滞考察的是筛选之后的路径如何进入现实。两者在代价问题上相接：计算与学习可以降低某些执行成本，也可能增加新的数据、验证和维护支出。改变是否值得，需要比较整条路径，而不是只计算候选出现的速度。

从比特、原子和生命三个介质看，这一点尤其清楚。比特空间中，路径搜索可以极快：文本、代码、图像、方案和模拟可以高速生成、修改、复制和撤回。原子空间中，路径选择必须落到材料、能量、位置、摩擦、热、力和维护上；错误不能像文本一样轻易撤销。生命空间中，路径选择更重，因为生命过程具有发育、代谢、免疫、繁殖和长期反馈；一个看似有利的干预，可能在更长时间尺度上显现代价。智能越深入这些介质，越需要承认：路径可以提前筛选，但状态转移仍要支付现实成本。

路径被选中之后，还要由材料、能源和执行结构实现。计算可以帮助减少这笔支出，却不能代替实际过程。下一步需要把资源条件展开，考察哪些支出来自初末状态，哪些来自工艺与损失。
